% !TeX root = ../papers/nc-p3.tex
\section{特異二次曲面像の共通頂点と反標準曲線の四次分岐}
\label{sec:singular-quadric-common-vertex}
\begin{theorem}\label{thm:singular-quadric-common-vertex}
十五次元閉包の有限四辺の正則整域窓で、四次中心系が基点を持ち、その像が三次元とする。
このとき $Z=k[x_0,\ldots,x_4,w]/(q,w\ell+f)$ の $q$ は必ず階数3で、
$q,\ell$ は共通頂点 $p$ を持ち、$f(p)\ne0$ である。$p$ からの射影で得る二次曲面錐の
最小解消 $S_0=\mathbb F_2$ 上で、$D=(\ell=0)$ の全引き戻しは被約な反標準曲線で、
滑らかな楕円曲線または節点を横断的につないだ有理曲線の巡回列（一つの節点有理曲線を含む）である。
一般の指数4の類は $k(S_0)$ の周期4の類の引き戻しで、$D$ の全成分で剰余の位数が4である。
例外曲線以外の各成分 $D_i$ の体上、頂点方向のモニック四次式 $f$ は一既約式の $m$ 乗で
$m=1,2,4$ である。その根体は剰余指標 $\chi_i$ の $m$ 倍が定める次数 $4/m$ の巡回拡大となる。
\end{theorem}
\begin{proof}[階数四の場合の頂点と因子類]
滑らかな $q$ は前節で除いた。まず階数4の $q$ を仮定する。
$Q=(q=0)$ の頂点の外は $S_0=\mathbb P^1\times\mathbb P^1$ 上の $L=\mathcal O(1,1)$ の全空間で、
$z$ をその繊維座標とする。$A=k[x]/q$ の類群の有理階数は1である。
実際四変数二行二列行列式の錐の類群は $\operatorname{Pic}(S_0)/\mathbb Z L$ であり、
多項式変数 $z$ の追加で変わらない。$Z[1/\ell]=A[1/\ell]$ であり、前々節の有限類群の定理から
この局所化の有理類群は零である。従って $\ell$ の成分の少なくとも一つが $A$ の非捩れ類を与える。

$\ell=a z^2+b z+c$、$a$ は定数、$b,c$ は $L,2L$ の切断である。
$a\ne0$ なら既約成分は全体または二つの一次グラフである。
モニック方程式の有理解は各正規チャートで正則なので一次グラフは $L$ の大域切断である。
成分はいずれも超平面の有理倍となり不可能である。
$a=0$, $b\ne0$ なら唯一の水平成分のほかは $b,c$ に共通な底の曲線だけである。
その全てが $L$ の倍数なら水平成分も同じになってしまう。
$(1,1)$ 切断 $b$ の真の非対角成分は罫線なので、例えば第一座標 $x$ の無限遠を成分に取れる。
もう一つの射影直線の体を $K=k(y)$ とすると、その一般点のアフィン平面上で極は
$P=b_0(y)z+c_1(y)x+c_0(y)=0$ の一直線である。ここで $b_0$ は第二座標の一次式である。
これは $b$ に既に第一座標の無限遠因子があることによる。
平面からこの直線を引いた部分で非分岐な類は $(\chi,P)$、$\chi\in H^1(K,\mathbb Z/4)$ である。
座標変更で $\mathbb G_m\times\mathbb A^1$ とし、零での剰余を取ればよい。$\operatorname{Br}(K)=0$ を使う。
$b_0$ の消えない任意の $y$ で $P$ は非定数一次式となり、この垂直面は極の成分でない。
そこでの記号の剰余から $\chi$ は非分岐である。座標の自明化で $P$ を $K$ の元で掛けても類は変わらない。
分岐候補が一次式の零の一か所だけでは射影直線に非自明な巡回指標はない。
従って $b=0$ でなければならず、$q,\ell$ は頂点を共有する。
\end{proof}

\begin{proof}[階数三の頂点方向]
階数3の二次三次元体の滑らかな部分は円錐の底 $\mathbb P^1$ 上の $\mathcal O(2)^{\oplus2}$ の全空間である。
繊維座標を $u,v$ とすると $\ell$ の最高二次部分は定数の二元二次式、一次係数は二次の底切断である。
最高部が非退化なら平方完成して $P=uv-d(t)$、$d$ は四次の切断とできる。
$d\ne0$ の場合、底のアフィン座標で被約既約な極の一般繊維は $\mathbb G_m$ であり、
その剰余は四乗類 $\delta(t)u^j$ の形である。有限な各 $t$ の上で $u$ が単元となる垂直曲線を選べる。
極のその一般点は滑らかであり、剰余の非分岐性から $\operatorname{ord}_t\delta$ は4の倍数である。
よって $\delta$ は四乗として落とせる。無限遠を $d$ の根でない位置に選ぶと $u$ の自明化は二位の極を持ち、
同じ検査で $2j=0\bmod4$ となる。また $d$ の重複度 $m$ の根で $v$ が単元の成分を取ると
$4\mid mj$ である。このためアフィン三空間上の記号 $(P,u)_4^j$ は $u=0$ では非分岐となり
（$d^j$ が四乗）、極上の指定した剰余を実現する。アフィン三空間のBrauer群は零なので
元の類もこの記号であり、$j$ が偶数だから指数は高々2である。
$d=0$ ならアフィン底上での補集合は $\mathbb A^1\times\mathbb G_m^2$、
その四捩れBrauer群は $(u,v)_4$ で生成される（Kummer列とトーラスのKünneth公式）。
無限遠で両座標の極次数が2なので、剰余の消滅は再び偶数倍を要求する。

最高部が階数1なら $P=u^2+b(t)v+c(t)$ とする。$b\ne0$ なら一般底体上で $P,u$ がアフィン座標で、
前段と同じく類は $(\chi,P)$ である。全ての底点で $P$ は非定数であり、
垂直面上のこの多項式は次数2以下、従って四乗でない。剰余の消滅は各底点で $\chi$ の
分岐位数が高々2であることを意味し、$2\chi$ は射影直線全体で非分岐、ゆえに零となる。
これは指数4を与えない。$b=0$ なら $v$ は既に共通頂点方向である。
最高部が零なら一般極は一次式で、二つの一次係数（底上の二次式）が共に消える点の外では
同じ議論で $\chi$ は非分岐である。係数が独立なら共通零点は一か所以下で不可能、
比例すれば一繊維方向が両式から落ちる。両方零でも共通方向がある。
\end{proof}

\begin{proof}[反標準分岐と頂点での四次係数]
以上のいずれの階数でも $q,\ell$ から一頂点変数 $z$ が落ちる。
残る二次曲面の滑らかな部分の一般点では、$z$ のアフィン直線の全体が $Q\setminus(\ell=0)$ にある。
従って $\alpha$ は底体の類 $\beta$ の引き戻しである。純超越拡大により指数は変わらず、
de Jongの定理より $\beta$ の周期も4である。底の滑らかなモデル $S_0$ は
階数4なら滑らかな二次曲面、階数3なら $\mathbb F_2$ であり、二次切断の引き戻し $D$ は反標準である。
$\beta$ は $D$ の外で非分岐である。滑らかな底の各因子については $z$ を不定元とした同じ剰余を使う。
$\mathbb F_2$ の負切断が $D$ に入らなければ $D$ と交わらず、そこに可能な剰余は
他の分岐曲線と会わない射影直線上の非分岐指標なのでやはり零である。

有理曲面上の非零な $2\beta$ の剰余の被約な支持 $C$ は $D$ の全体に等しい。
実際 $0<C<D$ なら $H^0(K_{S_0}+C)=0$ かつ $H^1(K_{S_0})=0$ なので随伴列から $H^1(\mathcal O_C)=0$。
このような被約曲線は正規化が有理で、成分と特異点を枝ごとにつないだグラフは森である
（正規化の完全列）。剰余の各枝の分岐数は、各成分上で和が零であり、各特異点でも
曲面剰余複体により和が零である。森の葉から消せば全部零となり、各指標も有理曲線上で非分岐となって零である。
滑らかな有理曲面のBrauer群は零なので矛盾する。よって $D$ は被約で全成分の $\beta$ の剰余が位数4である。
反標準曲線は連結で算術種数1。非零剰余には正規化の正種数かグラフ上の閉路が必要である。
前者ならその成分のみで種数1を使い切り、真の部分曲線の上記消滅から滑らかな楕円曲線だけとなる。
後者は真の部分曲線に閉路を持てず、算術種数1の全てがグラフの閉路によるので各枝は滑らかで
交点は通常節点である。平面曲線の三重点以上なら $\delta\ge r(r-1)/2>r-1$ ($r$ 枝) となるためである。

負切断を除く各 $D_i$ の一般点で $z$ に関する $f$ の特殊化を見る。
零多項式になると $w\ell+f$ が $\ell$ の一成分の上に余分な素成分を持ち、中心の整性に反する。
正次数の既約因子 $g$ が重複度 $m$ で現れれば、$w\ell+f=0$ における $D_i$ の対応する垂直因子は
分岐指数 $m$、剰余体は $k(D_i)[z]/g$ に独立な $w$ を添加したものである。
実際 $D_i$ は重複度1であり、その局所方程式を割った垂直因子一般点の長さはまさに $m$ である。
ここは正規中心の余次元1で次数1単元があるので $\alpha$ は非分岐で、
$k(D_i)[z]/g$ 上で $m\chi_i=0$ となる。従って $4\mid m\deg g$。

$f(p)=0$ なら $z$ の次数は高々3で、全ての $D_i$ で正次数の項が消えることになる。
従って $f$ の正の $z$ 係数は底の二次式 $\ell$ で割り切れる。実際被約なCartier因子の
全一般点で消え、正規な二次曲面上の切断を $\ell$ で割った有理切断は高さ1で正則である。
完全二次曲面の全切断は斉次座標環に一致する。$w$ の移動で $f$ から $z$ を除くと $Z=Z'[z]$。
この三次元の正規Gorenstein環 $Z'$ は $\omega_{Z'}=Z'$ で有理特異点ではあり得ない。
実際標数零の正次数正規有理特異点環の最高局所コホモロジーは零以上の次数で消える
（切断環の有理特異点の必要条件、または有理特異点の負の $a$ 不変量の定理）。
次数零で局所双対に反する。従って $f(p)\ne0$。各成分で四次モニックにでき、
先の整除から一既約因子だけで、根体は $m\chi_i$ の巡回拡大を次数一致によりちょうど与える。

最後に階数4を除く。$A$ 内で $\ell$ の異なる斉次素因子を $h$ 個とする。
それらの類が有理階数1を張るため、$A[1/\ell]$ の単元を定数で割った群の階数は $h-1$ である
（局所化の因子完全列）。一方 $f$ が頂点変数でモニックなので、各因子の底曲線上の
根による垂直素因子が $Z$ に一つ以上あり、互いに異なる。この計算をアフィン錐で行っても、
次数1の非零座標の局所自明化を加えれば同じである。$Z[1/\ell]=A[1/\ell]$ の単元だけでは
この $h$ 個以上の素因子の全ての有理類を消せず、$\operatorname{Cl}(Z)$ の有限性に反する。
\end{proof}
